\seksjon{5.3}{Sinussetningen}
\undertittel{Finn sider og vinkler i vilkårlige trekanter, og se når det finnes to løsninger.}
\maal{Begrunne og bruke sinussetningen, også når to trekanter passer med opplysningene}
\kompetansemaal{13,12}

\hvakanjeg{
\kan{1}{forklare hvorfor $\frac{\sin A}{a} = \frac{\sin B}{b} = \frac{\sin C}{c}$ ut fra arealsetningen}
\kan{1}{finne tredje vinkel med vinkelsummen}
\kan{1}{finne ukjente sider når to vinkler og én side er kjent}
\kanskille
\kan{2}{finne vinkel $B$ når to sider og en vinkel som ikke ligger mellom dem er kjent}
\kan{2}{avgjøre om det finnes ingen, én eller to trekanter, og sjekke vinkelsummen}
\kan{2}{forklare det tvetydige tilfellet med en svingarm}
}
